\section{Route implications and arithmetic boundary}\label{sec:boundary}

Theorem~\ref{thm:algebraic-exhaustion} is a genuine all-period Route-A
advance.  It supplies a complete primitive periodic-point layer, an exact
symbolic zeta on the real chain recurrent set, and totally real reflection
algebras.  In the repository evaluation vocabulary, the strongest tuple is
\[
(A1_{\rm analytic},A2_{\rm inherited\ symbolic},
 A3_{\rm partial},A4_{\rm formal}).
\]
The overall verdict remains \texttt{ROUTE\_A\_EXPLORATORY}.

The following statements are \emph{not} consequences of the theorem:
\begin{enumerate}
\item there is no intrinsic bijection between primitive H\'enon orbits and
rational primes;
\item no amplitude $\log p/p^{r/2}$ or von Mangoldt trace has been derived;
\item the full-shift zeta is not the completed Riemann $\xi$ function;
\item total reality of finite algebras is not self-adjointness of one
infinite-dimensional operator;
\item no Hilbert space, operator domain, compact resolvent, or zeta-regularized
determinant identity has been constructed.
\end{enumerate}
Route B is therefore not testable and is not authorized.

The next non-micro object is an all-period algebraic pressure.  If
$h(\alpha)$ is a canonical height or Galois-excess statistic on roots of
$\Psi_n$, one may ask for the existence and regularity of
\[
\mathcal P_{\rm Gal}(s)
=\limsup_{\substack{n\to\infty\\ n\,\mathrm{odd}}}
\frac1n\log\sum_{\Psi_n(\alpha)=0}e^{-s h(\alpha)}.
\]
A comparison with the physical instability pressure would be meaningful
without forcing a termwise prime labeling.  This is the narrow successor
problem retained by the present paper.
